% Implementation notes for the -sys and -lib common options.
%
% Standalone document: build with
%   latexmk -synctex=1 -pdf LibSysOptions.tex
%
% To fold it into the manual instead, delete everything up to \begin{document}
% (and the \maketitle/\tableofcontents lines), then \input it from ReadData.tex
% next to Structs and ReadSystemData; the \St, \Error, \Break, \Continue and
% \hltt definitions below are the ones the Vtf* pseudocode files already use,
% so they would move into the manual's preamble.
\documentclass[12pt,a4paper]{report}

% packages (same set as manual.tex, minus what this document does not use) %{{{
\usepackage[margin=3cm]{geometry}
\usepackage{amsmath}
\usepackage{array}
\usepackage{longtable}
\usepackage{booktabs}
\usepackage{etoolbox}
\usepackage[main=british,czech]{babel}
\usepackage[strings]{underscore}
\usepackage{xspace}
\usepackage{xcolor}
\usepackage{titleps}
\usepackage{tabularx}
\usepackage[shortlabels]{enumitem}
\usepackage[style=british]{csquotes}
\usepackage{verbatim}
\usepackage{listings}
\usepackage[normalem]{ulem}
\usepackage{setspace}
\usepackage{soul}
\usepackage{tikz}
% pseudocode
\usepackage{algpseudocode}
% hyperref should be the last package listed
\usepackage{hyperref}
\usepackage[nameinlink]{cleveref}
\usepackage{macros} %}}}

% pseudocode shorthands used by the Vtf* files, which never defined them %{{{
\algnewcommand\algorithmicbreak{\textbf{break}}
\algnewcommand\Break{\State \algorithmicbreak}
\algnewcommand\algorithmiccontinue{\textbf{continue}}
\algnewcommand\Continue{\State \algorithmiccontinue}
\algnewcommand\algorithmicerror{\textbf{error:}}
\algnewcommand\Error{\State \algorithmicerror\ }
\algnewcommand\algorithmicwarn{\textbf{warn:}}
\algnewcommand\Warn{\State \algorithmicwarn\ }
\let\St\State
% A \Comment longer than the space left on the line used to wrap to the left
% margin, so its tail sat under the code rather than under the comment. Set it
% in a block of its own instead, so the tail hangs under the text following the
% triangle. The block is only as wide as the comment needs, up to
% \algcommentmax, so short comments still sit flush right exactly as before;
% the \hfill-penalty-\hfill lets TeX drop the whole block to the next line when
% the code leaves too little room, rather than running into the margin.
\newlength{\algcommentwidth}
\newlength{\algcommenthang}
\newcommand{\algcommentmax}{0.45\linewidth}
\algrenewcommand\algorithmiccomment[1]{%
  \hfill\penalty0\hfill
  \settowidth{\algcommentwidth}{\(\triangleright\)~#1}%
  \ifdim\algcommentwidth>\algcommentmax\relax
    \setlength{\algcommentwidth}{\algcommentmax}%
  \fi
  \parbox[t]{\algcommentwidth}{%
    \settowidth{\algcommenthang}{\(\triangleright\)~}%
    \raggedright\hangindent=\algcommenthang\hangafter=1
    \(\triangleright\)~#1}%
}
\newcommand{\hltt}[1]{\hl{\tt{#1}}}
\newcommand{\algbottomrule}{\vspace{-0.5em}\noindent\rule{\textwidth}{0.4pt}} %}}}

\renewcommand{\mytitle}{Implementation of the -sys and -lib options}
\hypersetup{pdftitle={\mytitle}}

\begin{document}
\maketitle
\tableofcontents

\chapter{The \tt{-sys} and \tt{-lib} options}\label{chap:LibSys}

The two common options change what a system's types are \emph{called} rather
than what is in it: \tt{-sys <file>} names the molecule types and \tt{-lib
<dir>} names the bead types (and gives them their masses, charges, and radii).
These options make sense of a file that carries coordinates and topology but no
useful names, typically a \data file produced by LAMMPS.

The options are not independent: the library matches molecules \emph{by name},
so on a nameless input, \tt{-lib} can do its job only if \tt{-sys} has run
first. That ordering is fixed internally, the options' order in the command does
not matter.

\todo{Only \tt{-sys} and only \tt{-lib} comment?}

\section{Where they sit in a utility's start-up}\label{sec:LibSysOrder}

Every utility follows the following three steps (in that order):

\begin{enumerate}[nosep,leftmargin=20pt]
  \item \ttb{CommonOptions()} records the file and directory names
  \item \ttb{ReadStructure()} builds the \ttb{SYSTEM} described by the input
    file
  \item \ttb{ApplyLibraryOptions()} applies \tt{-sys}, then \tt{-lib}
\end{enumerate}

\noindent
Step 3 must come before any option that names a type (e.g., \tt{-bt},
\tt{-mt}, or a molecule-name argument), so that those options can use the
library's names rather than the input file's.

\section{Recording the options}\label{sec:LibSysRecord}

\ttb{CommonOptions()} stores both file and directory names in the
\ttb{COMMON_OPT} struct. An empty string marks \enquote{unused option}
throughout.

\begin{algorithmic}[1]
  \St \ttb{opt.lib} $\gets$ \tt{""}
  \St \ttb{FileOption}(\tt{-lib}, \ttb{opt.lib})\Comment{fills it in only if
    \tt{-lib} is present}
  \St \ttb{opt.sys} $\gets$ \tt{""}
  \St \ttb{FileOption}(\tt{-sys}, \ttb{opt.sys})
\end{algorithmic}
\algbottomrule

% \noindent
% Note that this happens for \emph{every} utility calling \ttb{CommonOptions()},
% whether or not it lists the options in its own help text.

\section{Applying the options}\label{sec:LibSysApply}

\CallInit{void ApplyLibraryOptions}{ReadLibrary.c}
\begin{longtable}{L{0.30\textwidth}Mp{0.545\textwidth}}
  \toprule
  parameter & in/out & explanation \\
  \midrule
  \ttb{(COMMON_OPT) commons} & in    & holds the \tt{-sys} and \tt{-lib} file
    and dir names \\
  \ttb{(SYSTEM) sys}         & in/out& the system whose types get renamed \\
  \ttb{(LIBRARY) lib}        & out   & the library dir; used only when
    \ttb{free_library} is \tt{false} \\
  \ttb{(bool) free_library}  & in    & whether the library is freed here
    (\tt{true} for all but \tt{Info}) \\
  \bottomrule
\end{longtable}

\begin{algorithmic}[1]
  \If{\ttb{commons.sys} $\neq$ \tt{""}}
    \St \ttb{ReadSysInfo}(\ttb{commons.sys}, \ttb{sys})
  \EndIf
  \If{\ttb{commons.lib} $\neq$ \tt{""}}
    \St \ttb{ApplyLibraryToSystem}(\ttb{commons.lib}, \ttb{sys}, \ttb{lib},
      \ttb{free_library})\ \Comment{\tt{-sys} has already run}
  \EndIf
\end{algorithmic}
\algbottomrule

\noindent
\ttb{free_library} decides who owns the library once the renaming is done, and
is the only thing separating the two kinds of caller:

\begin{itemize}[nosep,leftmargin=20pt]
  \item \tt{true}: the library is read, used and freed inside the call, and
    \ttb{lib} is ignored (case for most utilities)
  \item \tt{false}: the library is handed back in \ttb{lib}, so the caller
    must use \ttb{FreeLibrary()} to free it
\end{itemize}

\noindent
Only \tt{Info} passes \tt{false}, because it needs the library to print the DPD
interaction parameters.
% Note that \ttb{lib} is left untouched when \tt{-lib} was not given at all, so
% such a caller should initialise it and test \ttb{commons.lib} itself rather
% than expecting a flag back.

\section{\tt{-sys}: naming molecule types}\label{sec:SysOption}

\tt{Info} writes the \tt{-sys} file with \ttb{WriteSysInfo()}: one \tt{<mol
name> <count>} line per molecule type. Reading it back is purely
\emph{positional}: the $i$-th valid line names the $i$-th molecule type, meaning
the molecule types to rename in the input structure file must be in the same
order. The \tt{<count>} is parsed only to decide whether a line is a valid
entry; nothing cross-checks it against the system.

\CallInit{void ReadSysInfo}{ReadWriteSysInfo.c}
\begin{longtable}{L{0.30\textwidth}Mp{0.545\textwidth}}
  \toprule
  parameter & in/out & explanation \\
  \midrule
  \ttb{(char) filename[]} & in    & the \tt{-sys} file \\
  \ttb{(SYSTEM) System}   & in/out& molecule types get \ttb{Name} and
    \ttb{Named} \\
  \bottomrule
\end{longtable}

\begin{algorithmic}[1]
  \St open \ttb{filename}
  \St $i \gets 0$\Comment{index of the molecule type being named}
  \While{$line \gets$ next line of \ttb{filename}}
    \If{$line$ is empty, starts with \#, has fewer than 2 words, or 2nd word of
      $line$ is not a whole number}
      \Continue\Comment{skip blank, comment, and invalid lines}
    \EndIf
    \If{$i \geq$ \ttb{System.Count.MoleculeType}}
      \Warn more entries than molecule types; only the first $i$ are used
      \Break
    \EndIf
    \St \ttb{System.MoleculeType[$i$].Name} $\gets$ 1st word of $line$
    \St \ttb{System.MoleculeType[$i$].Named} $\gets$ \tt{true}
    \St $i \gets i + 1$
  \EndWhile
  \St close \ttb{filename}
  \If{$i <$ \ttb{System.Count.MoleculeType}}
    \Warn $i$ entries for a system with \ttb{Count.MoleculeType} molecule types
  \EndIf
\end{algorithmic}
\algbottomrule

\noindent
Two notes:

\begin{itemize}[nosep,leftmargin=20pt]
  \item because the match is positional, a file whose lines are in the wrong
    order will misname every type (both warnings above are about \emph{counts of
    entries}, not about whether the names make sense)
  \item \tt{-sys} touches no bead types; on a nameless \data file, it enables
    \tt{-lib} to work (\cref{sec:LibPass1}), but on its own (no \tt{-lib}
    option) it changes only the
    molecule names
\end{itemize}

\section{\tt{-lib}: naming bead types}\label{sec:LibOption}

\subsection{What the library holds}\label{ssec:LibContents}

\ttb{ReadLibrary()} reads four \tt{list_*.txt} tables from the library
directory into a \ttb{LIBRARY}:

\begin{longtable}{p{0.42\textwidth}p{0.515\textwidth}}
  \toprule
  file & what it contributes \\
  \midrule
  \tt{list_beadtypes.txt} & bead types (name, mass, charge, radius) and their
    self-interactions \\
  \tt{list_bonds.txt}      & bond types, each with a string ID \\
  \tt{list_angles.txt}     & angle types, each with a string ID \\
  \tt{list_cross_interactions.txt} & interactions between unlike bead types \\
  \bottomrule
\end{longtable}

\noindent
Everything else in the directory is one molecule per \tt{<mol name>.txt} file,
read on demand by \ttb{ReadLibraryMolFile()}. Those files are what the renaming
below keys on; the tables only supply the bead data that gets copied across.

\subsection{The top level}\label{ssec:LibApply}

\CallInit{void ApplyLibraryToSystem}{ReadLibrary.c}
\begin{longtable}{L{0.30\textwidth}Mp{0.545\textwidth}}
  \toprule
  parameter & in/out & explanation \\
  \midrule
  \ttb{(char) lib_dir[]}    & in    & the \tt{-lib} directory \\
  \ttb{(SYSTEM) sys}        & in/out& the system whose bead types get renamed \\
  \ttb{(LIBRARY) lib}       & out   & the library, for the caller to
    \ttb{FreeLibrary()}; ignored when \ttb{free_library} \\
  \ttb{(bool) free_library} & in    & whether the library is freed here \\
  \bottomrule
\end{longtable}

\begin{algorithmic}[1]
  \St $lib \gets$ \ttb{lib}, or a local variable if
    \ttb{free_library}\Comment{a null \ttb{lib} is an error otherwise}
  \St $before \gets$ copy of \ttb{sys.BeadType[]}\Comment{to tell afterwards
    whether anything happened}
  \St $lib \gets$ \ttb{ReadLibrary}(\ttb{lib_dir})
  \St \ttb{RenameBeadTypesFromLibrary}(\ttb{sys}, $lib$, \ttb{lib_dir})
  \St $changed \gets$ \tt{true} if any bead type's name, mass, charge or radius
    differs from $before$
  \St $unmatched \gets$ number of molecule types in \ttb{sys} with no
    \tt{<Name>.txt} in \ttb{lib_dir}
  \If{not $changed$}
    \Warn no bead type matched the library
  \ElsIf{$unmatched > 0$}
    \Warn $unmatched$ molecule type(s) have no library file, so the beads inside
      them keep what the input gave them; name them with \tt{-sys} if the input
      does not
  \EndIf
  \If{\ttb{free_library}}
    \St \ttb{FreeLibrary}($lib$)
  \EndIf
\end{algorithmic}
\algbottomrule

\noindent
A library that changes nothing is almost certainly the wrong directory, or an
input whose molecule types are unnamed, so it is worth a warning rather than
carrying on quietly with the file's own (possibly missing) masses.

\ttb{RenameBeadTypesFromLibrary()} then does the real work in two passes: bead
types used inside molecules are handled by molecule name, and free (unbonded)
bead types by charge. The two sets do not overlap, so neither pass can undo the
other.

\subsection{Pass 1: beads inside named molecules}\label{sec:LibPass1}

A molecule type whose name matches a library file lends its bead names,
position by position. The catch is that one \emph{system} bead type can be
reached from several \emph{library} beads at once --- a \data file with no
\tt{Masses} section has nothing to tell its neutral beads apart, so they all
arrive as a single type. Renaming as we go would give such a type whichever
library name came last, so proposals are collected first and applied only once
every molecule has been seen.

\begin{algorithmic}[1]
  \For{each bead type $i$ of \ttb{sys}}
    \St $proposed[i] \gets$ \emph{none}
  \EndFor
  \For{each molecule type $mt$ of \ttb{sys}}
    \If{\ttb{$mt$.Name} is empty}
      \Continue
    \EndIf
    \If{the file named \ttb{$mt$.Name}\ttb{.txt} cannot be read from
        \ttb{lib_dir}}
      \Continue
    \EndIf
    \St $mol \gets$ that file
    \If{$mol$.n\_beads $\neq$ \ttb{$mt$.nBeads}}
      \Continue\Comment{not the molecule the name suggests; counterions are
        separate molecules, so the counts must agree}
    \EndIf
    \For{each bead position $b$ of $mol$}
      \St $st \gets$ \ttb{$mt$.Bead[$b$]}\Comment{system bead type at that
        position}
      \St $lt \gets$ \ttb{FindBeadType}($mol$.bead\_name[$b$],
        \ttb{lib})\Comment{by name}
      \If{$lt$ is \emph{none}}
        \Continue
      \EndIf
      \If{$proposed[st]$ is \emph{none}}
        \St $proposed[st] \gets lt$
      \ElsIf{$proposed[st] \neq lt$}
        \St $proposed[st] \gets$ \emph{ambiguous}
      \EndIf
    \EndFor
  \EndFor
  \St $ambiguous \gets 0$
  \For{each bead type $i$ of \ttb{sys}}
    \If{$proposed[i]$ is \emph{ambiguous}}
      \St $ambiguous \gets ambiguous + 1$
      \Continue\Comment{left exactly as the input had it}
    \EndIf
    \If{$proposed[i]$ is \emph{none}}
      \Continue
    \EndIf
    \St copy \ttb{Name}, \ttb{Charge}, \ttb{Mass} and \ttb{Radius} from library
      bead type $proposed[i]$
  \EndFor
  \If{$ambiguous > 0$}
    \Warn $ambiguous$ bead type(s) stand for several library beads at once and
      are left alone; the input cannot tell them apart
  \EndIf
\end{algorithmic}
\algbottomrule

\noindent
The guiding rule, here and in pass 2: \emph{a wrong name is worse than no
name}.

\subsection{Pass 2: free beads}\label{sec:LibPass2}

Beads that belong to no molecule cannot be reached by name, because there is no
molecule name to look up. They are matched by charge instead, and by count where
the count is predictable. Candidates come from walking the library directory:

\begin{algorithmic}[1]
  \For{each bead type $bt$ used by any molecule type of \ttb{sys}}
    \St $in\_mol[bt] \gets$ \tt{true}\Comment{pass 1's territory; skipped below}
  \EndFor
  \St $cand \gets$ empty list
  \For{each \tt{*.txt} in \ttb{lib_dir} that is not a \tt{list_*} table}
    \St $mol \gets$ that file, skipping it if unreadable
    \If{$mol$ has exactly one bead, known to the library}
      \St append $\{$name, charge, $count = $ \emph{unknown}$\}$ to
        $cand$\Comment{any single-bead molecule is loose in the system, salts
        included}
    \EndIf
    \St $parent \gets$ number of molecules in \ttb{sys} whose type is named
      $mol$.name
    \If{$parent = 0$}
      \Continue\Comment{its counterions are not in this system}
    \EndIf
    \For{each counterion $c$ declared by $mol$}
      \St $cion \gets$ the file named \ttb{$c$.name}\ttb{.txt}, skipping it if
        unreadable or not monoatomic
      \St find or append the $cand$ entry for $c$.name, with $count$ starting
        at 0
      \St $count \gets count + parent \times c$.count\Comment{summed over every
        parent molecule}
    \EndFor
  \EndFor
\end{algorithmic}
\algbottomrule

\noindent
Matching then goes from the most reliable evidence to the least. A name the
library already knows beats any guess made from charge --- the file has said
what the type is, so there is nothing to work out:

\begin{algorithmic}[1]
  \For{each bead type $bt$ with $in\_mol[bt]$ false}
    \St $lt \gets$ \ttb{FindBeadType}(\ttb{sys.BeadType[$bt$].Name}, \ttb{lib})
    \If{$lt$ is \emph{none}}
      \Continue
    \EndIf
    \St copy \ttb{Charge}, \ttb{Mass} and \ttb{Radius} from library bead type
      $lt$\Comment{the name already \emph{is} the library's}
    \St $in\_mol[bt] \gets$ \tt{true}\Comment{done with; the charge matching
      below skips it}
  \EndFor
  \For{each bead type $bt$ with $in\_mol[bt]$ still false}
    \St $q \gets$ \ttb{sys.BeadType[$bt$].Charge},\quad $n \gets$
      \ttb{sys.BeadType[$bt$].Number}
    \St $match \gets$ the \emph{only} $cand$ with a known $count$,
      $|charge - q| < 0.01$ and $count = n$
    \If{$match$ is \emph{none}}
      \St $match \gets$ the \emph{only} $cand$ with an unknown $count$ and
        $|charge - q| < 0.01$
    \EndIf
    \If{$match$ is not \emph{none}}
      \St \ttb{Name} $\gets match$.name; copy \ttb{Charge}, \ttb{Mass} and
        \ttb{Radius} from the library
    \EndIf
  \EndFor
\end{algorithmic}
\algbottomrule

\noindent
Both \enquote{the only} steps rename nothing when two or more candidates fit.
Note that an ambiguous charge-plus-count match does not give up: it falls
through to the charge-only step, which may still find a unique answer among the
candidates whose count is unpredictable.

Leaving single-bead molecules that declare counterions of their own out of the
candidate list is what once made an \tt{Na+} bead get renamed to the
identically charged \tt{H+}; they are included precisely because they have no
predictable count and can only be matched on charge.

\subsection{Summary of the warnings}\label{ssec:LibWarnings}

\begin{longtable}{>{\raggedright\arraybackslash}p{0.44\textwidth}p{0.495\textwidth}}
  \toprule
  warning & what it means \\
  \midrule
  \tt{no bead type matched the library} & nothing at all changed: wrong
    directory, or an input whose molecule types are unnamed and whose free
    beads match nothing \\
  \tt{<n> molecule type(s) have no library file} & those molecules' beads keep
    what the input gave them --- add \tt{-sys} if the input does not name its
    molecules \\
  \tt{<n> bead type(s) stand for several library beads} & pass 1 found
    conflicting proposals; the input cannot tell those beads apart, so they are
    left alone \\
  \bottomrule
\end{longtable}

\section{Why the order matters}\label{sec:LibSysOrder2}

\tt{Examples/SystemID/01-FilesAndOptions} builds the same system in several
file formats and then strips information out of them, which makes the division
of labour easy to see. Starting from \tt{foreign.data} --- a \data file with no
name comments, so bead types arrive as \tt{b0}\ldots\tt{b4} and the molecule
type as \tt{m}:

\begin{longtable}{p{0.30\textwidth}p{0.635\textwidth}}
  \toprule
  options & result \\
  \midrule
  \tt{-sys} alone & molecule types get their names; every bead type is
    untouched \\
  \tt{-lib} alone & pass 1 finds no molecule file to match, so only the free
    beads are identified (water by charge, the counterion by charge and count);
    warns \\
  \tt{-lib} \emph{and} \tt{-sys} & \tt{-sys} runs first, so pass 1 can match
    the molecule and every bead type is identified \\
  \bottomrule
\end{longtable}

\noindent
The refusal case is worth seeing too: on \tt{nomass.data} (the same file with
the \tt{Masses} section removed) even \tt{-lib} with \tt{-sys} leaves the
molecule's neutral beads alone, because they arrive as one bead type standing
for three library beads --- pass 1's \emph{ambiguous} outcome.

\section{Utilities that do more}\label{sec:LibSysExtra}

\begin{itemize}[nosep,leftmargin=20pt]
  \item \tt{Info} declares its own \tt{-lib} and \tt{-sys} help entries and is
    the one caller passing \ttb{free_library} as \tt{false}, to print the DPD
    interaction parameters and to write a \field file's interactions. It also has
    \tt{--check-library} and \tt{--mol-info}, which need only \tt{-lib} and no
    input file at all.
  \item \tt{AddToSystem} uses \tt{-lib} as a \emph{source of molecules} to
    add via \tt{-mol}, which is a different job from renaming what is already
    there. Its \tt{-sys} takes an optional second file name to write the
    updated composition to, and when building a system from scratch a single
    argument is treated as output rather than input.
\end{itemize}

\end{document}
